\setcounter{chapter}{4}
\pagestyle{fancy}

\chapter{线性变换}

\section{线性变换及其矩阵表示}

\begin{theorem}[\markstar]
设 $\sigma$ 是 $n$ 维线性空间 $V$ 的线性变换,则下述命题等价:
\begin{enumerate}
    \item $\sigma$ 是可逆线性变换;
    \item $\sigma$ 是单射;
    \item $\sigma$ 是满射;
    \item 若 $\eta_1,\eta_2,\ldots,\eta_n$ 是 $V$ 的一个基,则
    $\sigma(\eta_1),\sigma(\eta_2),\ldots,\sigma(\eta_n)$ 也是 $V$ 的一个基;
    \item 若 $W,U$ 是 $V$ 的两个子空间,且 $V=W\oplus U$,则
    \[
        V=\sigma(W)\oplus\sigma(U).
    \]
\end{enumerate}
\end{theorem}

\begin{example}
设 $\mathbb{R}^3$ 的线性变换 $\sigma$ 在基
\[
\alpha_1=(1,1,0)^{\mathsf T},\quad
\alpha_2=(2,-1,0)^{\mathsf T},\quad
\alpha_3=(0,2,-1)^{\mathsf T}
\]
下的矩阵为
\[
A=
\begin{pmatrix}
2&0&3\\
0&-2&-1\\
1&-1&4
\end{pmatrix}.
\]
向量 $\alpha$ 在基
\[
\beta_1=(1,1,3)^{\mathsf T},\quad
\beta_2=(1,2,5)^{\mathsf T},\quad
\beta_3=(0,2,-1)^{\mathsf T}
\]
下的坐标为 $(1,0,-1)^{\mathsf T}$,求 $\sigma(\alpha)$ 在基
$\beta_1,\beta_2,\beta_3$ 下的坐标.
\end{example}
\exampleblank

\begin{solution}
由
\[
\alpha=\beta_1-\beta_3=(1,-1,4)^{\mathsf T}.
\]
设
\[
P=(\alpha_1,\alpha_2,\alpha_3),\qquad
Q=(\beta_1,\beta_2,\beta_3).
\]
则
\[
[\alpha]_{\alpha}=P^{-1}\alpha
=
\begin{pmatrix}
5\\-2\\-4
\end{pmatrix},
\qquad
[\sigma(\alpha)]_{\alpha}
=
A[\alpha]_{\alpha}
=
\begin{pmatrix}
-2\\8\\-9
\end{pmatrix}.
\]
故
\[
\sigma(\alpha)
=
P
\begin{pmatrix}
-2\\8\\-9
\end{pmatrix}
=
\begin{pmatrix}
14\\-28\\9
\end{pmatrix}.
\]
再换到 $\beta$ 基下:
\[
[\sigma(\alpha)]_{\beta}
=
Q^{-1}\sigma(\alpha)
=
\boxed{
\begin{pmatrix}
\dfrac{178}{5}\\[2pt]
-\dfrac{108}{5}\\[2pt]
-\dfrac{51}{5}
\end{pmatrix}
}.
\]
\end{solution}


\begin{example}
设 $P[x]_2$ 表示实数域上次数不超过 $2$ 的多项式构成的线性空间,
\[
f_1=1-x,\qquad f_2=1+x^2,\qquad f_3=x+2x^2
\]
是 $P[x]_2$ 的一个基.线性变换 $\sigma$ 满足
\[
\sigma(f_1)=2+x^2,\qquad
\sigma(f_2)=x,\qquad
\sigma(f_3)=1+x+x^2.
\]
\begin{enumerate}
    \item 求由基 $1,x,x^2$ 到基 $f_1,f_2,f_3$ 的过渡矩阵;
    \item 求 $\sigma$ 在基 $f_1,f_2,f_3$ 下的矩阵;
    \item 设 $f=1+2x+3x^2$,求 $\sigma(f)$.
\end{enumerate}
\end{example}
\exampleblank

\begin{solution}
在标准基 $1,x,x^2$ 下,
\[
F=
\begin{pmatrix}
1&1&0\\
-1&0&1\\
0&1&2
\end{pmatrix},
\qquad
(f_1,f_2,f_3)=(1,x,x^2)F.
\]
因此由标准基到 $f_1,f_2,f_3$ 的过渡矩阵为
\[
\boxed{
F^{-1}
=
\begin{pmatrix}
-1&-2&1\\
2&2&-1\\
-1&-1&1
\end{pmatrix}
}.
\]

又
\[
S=
\begin{pmatrix}
2&0&1\\
0&1&1\\
1&0&1
\end{pmatrix}
\]
的列分别是 $\sigma(f_i)$ 在标准基下的坐标,故
\[
[\sigma]_{(f)}
=
F^{-1}S
=
\boxed{
\begin{pmatrix}
-1&-2&-2\\
3&2&3\\
-1&-1&-1
\end{pmatrix}
}.
\]
最后
\[
[f]_{(f)}=F^{-1}
\begin{pmatrix}1\\2\\3\end{pmatrix},
\]
计算得
\[
\boxed{\sigma(f)=-4+3x-2x^2}.
\]
\end{solution}


\begin{example}
已知 $\mathbb{R}^3$ 的线性变换 $\sigma$ 对基
\[
\varepsilon_1=(-1,0,2)^{\mathsf T},\quad
\varepsilon_2=(0,1,1)^{\mathsf T},\quad
\varepsilon_3=(3,-1,-6)^{\mathsf T}
\]
的像分别为
\[
\sigma(\varepsilon_1)=(-1,0,1)^{\mathsf T},\quad
\sigma(\varepsilon_2)=(0,-1,2)^{\mathsf T},\quad
\sigma(\varepsilon_3)=(-1,-1,3)^{\mathsf T}.
\]
\begin{enumerate}
    \item 求 $\sigma$ 在基 $\varepsilon_1,\varepsilon_2,\varepsilon_3$ 下的矩阵;
    \item 设 $x=(1,1,1)^{\mathsf T}$,求 $\sigma(x)$;
    \item 已知 $\sigma(x)$ 在基 $\varepsilon_1,\varepsilon_2,\varepsilon_3$ 下的坐标向量为 $(2,-4,2)^{\mathsf T}$,求 $x$;
    \item 证明 $\varepsilon_1,\varepsilon_1+\varepsilon_2,
    \varepsilon_1+\varepsilon_2+\varepsilon_3$ 是 $\mathbb{R}^3$ 的基,并求
    $\sigma$ 在该基下的矩阵.
\end{enumerate}
\end{example}
\exampleblank

\begin{solution}
记
\[
E=(\varepsilon_1,\varepsilon_2,\varepsilon_3),
\qquad
S=(\sigma(\varepsilon_1),\sigma(\varepsilon_2),\sigma(\varepsilon_3)).
\]
则
\[
[\sigma]_{\varepsilon}
=
E^{-1}S
=
\boxed{
\begin{pmatrix}
-2&9&7\\
-1&2&1\\
-1&3&2
\end{pmatrix}
}.
\]
由标准坐标矩阵
\[
[\sigma]_{\mathrm{std}}=SE^{-1}
\]
可得
\[
\boxed{
\sigma(1,1,1)^{\mathsf T}
=
(-7,-5,17)^{\mathsf T}
}.
\]

第三问按原题数据无解:将目标坐标换回标准坐标后,线性方程
$\sigma(x)=E(2,-4,2)^{\mathsf T}$ 不相容.故该小问题面或数据存在疑点,
已在《解答问题记录》中单独标记.

对第四问,令
\[
Q=
(\varepsilon_1,\,
\varepsilon_1+\varepsilon_2,\,
\varepsilon_1+\varepsilon_2+\varepsilon_3).
\]
从 $\varepsilon$ 基到该基的变换矩阵为上三角且行列式为 $1$,
故这三个向量构成一组基.计算得
\[
\boxed{
[\sigma]_Q
=
\begin{pmatrix}
-1&6&12\\
0&-1&-2\\
-1&2&4
\end{pmatrix}
}.
\]
\end{solution}


\begin{example}
设 $M_2(\mathbb{R})$ 中
\[
\eta_1=
\begin{pmatrix}1&0\\0&0\end{pmatrix},\quad
\eta_2=
\begin{pmatrix}1&1\\0&0\end{pmatrix},\quad
\eta_3=
\begin{pmatrix}1&1\\1&0\end{pmatrix},\quad
\eta_4=
\begin{pmatrix}1&1\\1&1\end{pmatrix}
\]
构成一组基,又设
\[
\xi_1=
\begin{pmatrix}1&0\\3&0\end{pmatrix},\quad
\xi_2=
\begin{pmatrix}1&1\\3&3\end{pmatrix},\quad
\xi_3=
\begin{pmatrix}3&1\\7&3\end{pmatrix},\quad
\xi_4=
\begin{pmatrix}3&3\\7&7\end{pmatrix}.
\]
若线性变换 $\sigma$ 满足 $\sigma(\eta_i)=\xi_i$,求:
\begin{enumerate}
    \item $\sigma(\xi_1),\ldots,\sigma(\xi_4)$;
    \item $\sigma(\xi_1),\ldots,\sigma(\xi_4)$ 能否构成 $M_2(\mathbb{R})$ 的一组基.
\end{enumerate}
\end{example}
\exampleblank

\begin{solution}
将二阶矩阵按 $(a_{11},a_{12},a_{21},a_{22})^{\mathsf T}$ 向量化.
计算得到
\[
\begin{aligned}
\sigma(\xi_1)&=
\begin{pmatrix}7&0\\15&0\end{pmatrix},&
\sigma(\xi_2)&=
\begin{pmatrix}7&7\\15&15\end{pmatrix},\\
\sigma(\xi_3)&=
\begin{pmatrix}17&7\\37&15\end{pmatrix},&
\sigma(\xi_4)&=
\begin{pmatrix}17&17\\37&37\end{pmatrix}.
\end{aligned}
\]
又 $\xi_1,\ldots,\xi_4$ 的坐标矩阵行列式非零,而 $\sigma$ 在
$M_2(\mathbb{R})$ 上的矩阵行列式也非零,因此 $\sigma$ 可逆.
可逆线性变换保持线性无关性,故
\[
\boxed{\sigma(\xi_1),\sigma(\xi_2),\sigma(\xi_3),\sigma(\xi_4)
\text{ 构成一组基}.}
\]
\end{solution}


\begin{example}
设 $V$ 是数域 $P$ 上的 $n$ 维线性空间,
$\alpha_1,\ldots,\alpha_n$ 为一组基.定义线性变换 $A$:
\[
A(\alpha_i)=\alpha_{i+1}\quad(i=1,\ldots,n-1),
\qquad
A(\alpha_n)=0.
\]
\begin{enumerate}
    \item 求 $A$ 在该基下的矩阵;
    \item 证明 $A^n=O$,但 $A^{n-1}\neq O$;
    \item 若 $B$ 为线性变换且 $B^n=O,B^{n-1}\neq O$,证明存在一组基使 $B$ 的矩阵也等于上述矩阵.
\end{enumerate}
\end{example}
\exampleblank

\begin{solution}
矩阵为
\[
\boxed{
\begin{pmatrix}
0&0&\cdots&0\\
1&0&\cdots&0\\
0&1&\ddots&\vdots\\
\vdots&\ddots&\ddots&0\\
0&\cdots&1&0
\end{pmatrix}
}.
\]
由定义,
\[
A^k(\alpha_1)=\alpha_{k+1}\quad(k<n),
\]
故 $A^{n-1}\alpha_1=\alpha_n\neq0$,而 $A^n$ 对全部基向量都为零.

对第三问,由 $B^{n-1}\neq O$,可取 $v$ 使 $B^{n-1}v\neq0$.
于是
\[
B^{n-1}v,\ B^{n-2}v,\ \ldots,\ Bv,\ v
\]
线性无关;否则取最低次非零关系并作用适当次幂的 $B$ 即得矛盾.
以
\[
\beta_1=B^{n-1}v,\ \beta_2=B^{n-2}v,\ldots,\beta_n=v
\]
为基,即得到同一幂零 Jordan 块矩阵.
\end{solution}


\begin{example}
设 $f:\mathbb{R}^2\to\mathbb{R}$ 为线性映射.
\begin{enumerate}
    \item 证明存在 $a,b\in\mathbb{R}$ 使
    \[
    f(x,y)=ax+by;
    \]
    \item 已知 $f(1,1)=3,\ f(1,0)=4$,求 $f(2,1)$.
\end{enumerate}
\end{example}
\exampleblank

\begin{solution}
由于
\[
(x,y)=x(1,0)+y(0,1),
\]
故
\[
f(x,y)=xf(1,0)+yf(0,1).
\]
令 $a=f(1,0),b=f(0,1)$ 即得第一问.

第二问中 $a=4$,又 $a+b=3$,故 $b=-1$,
\[
\boxed{f(2,1)=7}.
\]
\end{solution}


\begin{example}[\markstar]
设 $A$ 为数域 $P$ 上的 $n$ 阶方阵,定义 $P^{n\times n}$ 上的线性变换
\[
T(X)=AX.
\]
求 $\operatorname{tr}T$ 与 $\det T$.
\end{example}
\exampleblank

\begin{solution}
按列分解矩阵空间:
\[
M_n(P)\cong P^n\oplus\cdots\oplus P^n
\]
共 $n$ 份,而 $T$ 在每一份上都由 $A$ 作用.因此 $T$ 的矩阵相似于
$n$ 个 $A$ 的块对角直和,故
\[
\boxed{
\operatorname{tr}T=n\,\operatorname{tr}A,
\qquad
\det T=(\det A)^n.
}
\]
\end{solution}


\begin{example}
设 $f$ 是平面 $\mathbb{R}^2$ 上的线性变换,满足:
\begin{enumerate}
    \item $f(1,0)$ 位于第四象限;
    \item $f(0,1)$ 位于第二象限;
    \item $f(1,1)$ 位于第一象限.
\end{enumerate}
证明 $f$ 可逆,且 $f^{-1}$ 把第一象限内任意点映射到第一象限.
\end{example}
\exampleblank

\begin{solution}
设
\[
f(1,0)=(a,-b),\qquad
f(0,1)=(-c,d),
\]
其中 $a,b,c,d>0$.由 $f(1,1)$ 位于第一象限知
\[
a>c,\qquad d>b.
\]
故
\[
\det
\begin{pmatrix}
a&-c\\
-b&d
\end{pmatrix}
=ad-bc>0,
\]
所以 $f$ 可逆.

若 $(u,v)$ 在第一象限,则
\[
f^{-1}
\begin{pmatrix}u\\v\end{pmatrix}
=
\dfrac{1}{ad-bc}
\begin{pmatrix}
d&c\\
b&a
\end{pmatrix}
\begin{pmatrix}u\\v\end{pmatrix},
\]
两个分量均为正,故结论成立.
\end{solution}


\begin{example}
设 $V$ 为数域 $F$ 上的 $n$ 维线性空间.
\begin{enumerate}
    \item 证明 $V$ 的任意有限多个真子空间的并仍是 $V$ 的真子集;
    \item 设 $A_1,\ldots,A_s$ 是 $V$ 上两两不同的线性变换,证明存在
    $v\in V$ 使 $A_i(v)$ 两两不同.
\end{enumerate}
\end{example}
\exampleblank

\begin{solution}
第一问即有限不覆盖定理.

第二问对每个 $i\neq j$,
\[
\ker(A_i-A_j)
\]
是真子空间,因为 $A_i\neq A_j$.有限不覆盖定理保证存在
\[
v\notin\bigcup_{i<j}\ker(A_i-A_j).
\]
于是
\[
(A_i-A_j)v\neq0,
\]
故 $A_i(v)\neq A_j(v)$.
\end{solution}


\begin{example}[\markstar]
设 $V,V'$ 是同一数域 $P$ 上的两个线性空间,$T:V\to V'$ 为线性映射,
并存在映射 $\widetilde T:V'\to V$.原题给出某一侧复合等于恒等映射,
讨论 $T,\widetilde T$ 的单射,满射与可逆性.
\end{example}
\exampleblank

\begin{solution}
这一题原 PDF 的复合顺序与恒等映射所在空间在类型上不一致.
正确的标准结论是:
\[
\widetilde T T=I_V
\quad\Longrightarrow\quad
T\text{ 单射, }\widetilde T\text{ 满射};
\]
而
\[
T\widetilde T=I_{V'}
\quad\Longrightarrow\quad
T\text{ 满射, }\widetilde T\text{ 单射}.
\]
单有一侧逆一般不能推出 $T$ 可逆;有限维且 $\dim V=\dim V'$ 时,
单射与满射等价,此时才可进一步推出可逆.
\end{solution}


\begin{example}[\markstar]
问是否存在 $n$ 阶方阵 $A,B$ 满足
\[
AB-BA=I_n?
\]
又问线性空间上的线性变换是否可能满足同样关系,并讨论何种条件下存在.
\end{example}
\exampleblank

\begin{solution}
在特征为 $0$ 的数域上不可能,因为
\[
\operatorname{tr}(AB-BA)=0,
\qquad
\operatorname{tr}I_n=n\neq0.
\]

更一般地,若底域特征为 $p>0$,必要条件是
\[
p\mid n.
\]
该条件也是充分的.例如在 $p$ 维空间
\[
P_{p-1}(F)=\Span\{1,x,\ldots,x^{p-1}\}
\]
上,令 $D=\dfrac{\mathrm d}{\mathrm dx}$,$X$ 为乘以 $x$ 后模 $x^p$ 的算子,
则在特征 $p$ 下
\[
DX-XD=I.
\]
当 $p\mid n$ 时作若干个这样的直和即可.
\end{solution}


\begin{example}
设 $V$ 是复数域 $\mathbb{C}$ 上的 $n$ 维线性空间,$\sigma$ 是 $V$ 的线性变换.
把同一加法与实数数乘保留下来后得到实线性空间 $V_{\mathbb{R}}$.
若 $\sigma$ 在复基下的矩阵为 $A_{\mathbb{C}}$,在某实基下的矩阵为
$A_{\mathbb{R}}$,证明
\[
\det A_{\mathbb{R}}
=
|\det A_{\mathbb{C}}|^2.
\]
\end{example}
\exampleblank

\begin{solution}
在由复基
\[
e_1,\ldots,e_n
\]
诱导的实基
\[
e_1,ie_1,\ldots,e_n,ie_n
\]
下,若
\[
A_{\mathbb{C}}=B+iC,
\]
则对应实矩阵为
\[
\begin{pmatrix}
B&-C\\
C&B
\end{pmatrix}.
\]
其行列式等于
\[
\det(B+iC)\det(B-iC)
=
\det A_{\mathbb{C}}\,
\overline{\det A_{\mathbb{C}}}
=
|\det A_{\mathbb{C}}|^2.
\]
\end{solution}


\begin{example}
设 $P$ 为数域,$\tau$ 为 $M_n(P)$ 上的线性变换.对任意固定
$A,B\in M_n(P)$,至少有
\[
\tau(AB)=\tau(A)\tau(B)
\]
或
\[
\tau(AB)=\tau(B)\tau(A)
\]
之一成立.证明:要么对所有 $A,B$ 总是前一种关系,要么对所有 $A,B$
总是后一种关系.
\end{example}
\exampleblank

\begin{solution}
令
\[
F(A,B)=\tau(AB)-\tau(A)\tau(B),
\qquad
G(A,B)=\tau(AB)-\tau(B)\tau(A).
\]
假设存在一对 $(A_0,B_0)$ 使 $F(A_0,B_0)\neq0$,则必有
$G(A_0,B_0)=0$.对任意 $A,B$,考虑参数族
\[
(A_0+tA,\ B_0+sB).
\]
题设保证对所有参数点总有 $F=0$ 或 $G=0$.由于 $F,G$ 的每个矩阵元
都是 $s,t$ 的多项式,而数域无限时有限个真代数零集不能覆盖整个平面,
从而其中一种恒为零.由在 $(A_0,B_0)$ 处 $F\neq0$,只能是 $G\equiv0$.
另一种情形同理.

若底域是有限域,上述论证需要额外处理;原题未说明数域是否默认无限,
这一点在问题记录中注明.
\end{solution}


\begin{lemma}[\markstar]
\begin{enumerate}
    \item 若 $F$ 上方阵 $A$ 有特征值 $\lambda_1\in F$,则存在 $F$ 上可逆矩阵 $P$ 使
    \[
    P^{-1}AP=
    \begin{pmatrix}
    \lambda_1&*\\
    0&B_1
    \end{pmatrix}.
    \]
    \item 若实方阵 $A$ 有非实特征值 $\lambda_1=a+bi$,相应复特征向量
    $x=y+zi$,则 $L(y,z)$ 是二维实不变子空间,并存在实可逆矩阵 $P$ 使
    \[
    P^{-1}AP=
    \begin{pmatrix}
    a&b&*\\
    -b&a&*\\
    0&0&C_1
    \end{pmatrix}.
    \]
\end{enumerate}
\end{lemma}


\begin{theorem}[\markstar]
设 $A$ 为域 $F$ 上的 $n$ 阶方阵.
\begin{enumerate}
    \item 若已知 $A$ 在 $F$ 中有特征值 $\lambda_1,\ldots,\lambda_t$,
    则存在相似变换把这些特征值依次放到上三角块的对角线上;
    \item 若 $A$ 在 $F$ 中有 $n$ 个特征值(重根按重数计算),则 $A$ 相似于上三角矩阵;
    \item 实方阵总能在实数域上相似于准上三角矩阵,其对角块为一阶实块或二阶实块.
\end{enumerate}
\end{theorem}


\begin{proposition}[\markstar]
设 $A,B$ 分别为 $n,m$ 阶矩阵,特征值为
$\lambda_i$ 与 $\mu_j$,则 $A\otimes B$ 的特征值为
\[
\lambda_i\mu_j.
\]
\end{proposition}


\begin{proposition}[\marktwostar]
设 $A,B$ 分别为数域 $F$ 上的 $m,n$ 阶方阵,
$V=M_{m\times n}(F)$,定义
\[
\varphi(X)=AXB.
\]
则:
\begin{enumerate}
    \item 在矩阵单位基下,$\varphi$ 的矩阵为 $A\otimes B^{\mathsf T}$;
    \item $\varphi$ 的特征值为 $\lambda_i\mu_j$.
\end{enumerate}
\end{proposition}


\begin{example}
设 $A,B$ 分别为数域 $F$ 上的 $m,n$ 阶方阵,
$V=M_{m\times n}(F)$,定义
\[
\varphi(X)=AX-XB.
\]
证明 $\varphi$ 的特征值为
\[
\lambda_i-\mu_j.
\]
\end{example}
\exampleblank

\begin{solution}
向量化公式给出
\[
\operatorname{vec}(AX-XB)
=
\left(I_n\otimes A-B^{\mathsf T}\otimes I_m\right)\operatorname{vec}X.
\]
前两项可同时上三角化,故该 Kronecker 和的对角元正是
\[
\lambda_i-\mu_j.
\]
因此这些数构成 $\varphi$ 的全部特征值.
\end{solution}


\begin{example}
设 $A=(a_{ij})$ 为 $n$ 阶复方阵,
\[
g(\lambda)=|\lambda I_n+A|.
\]
证明 $n^2$ 阶块矩阵
\[
B=
\begin{pmatrix}
a_{11}I_n+A&a_{12}I_n&\cdots&a_{1n}I_n\\
a_{21}I_n&a_{22}I_n+A&\cdots&a_{2n}I_n\\
\vdots&\vdots&\ddots&\vdots\\
a_{n1}I_n&a_{n2}I_n&\cdots&a_{nn}I_n+A
\end{pmatrix}
\]
可逆的充分必要条件是 $g(A)$ 可逆.
\end{example}
\exampleblank

\begin{solution}
注意
\[
B=A\otimes I_n+I_n\otimes A.
\]
若 $A$ 的特征值为 $\lambda_1,\ldots,\lambda_n$,则 $B$ 的特征值为
\[
\lambda_i+\lambda_j\qquad(1\leq i,j\leq n).
\]
故
\[
B\text{ 可逆}
\iff
\lambda_i+\lambda_j\neq0\quad\forall i,j.
\]
另一方面 $g(A)$ 的特征值为
\[
g(\lambda_i)=\prod\limits_{j=1}^{n}(\lambda_i+\lambda_j),
\]
故
\[
g(A)\text{ 可逆}
\iff
g(\lambda_i)\neq0\quad\forall i
\iff
\lambda_i+\lambda_j\neq0\quad\forall i,j.
\]
两者等价.
\end{solution}


\newpage

\section{值域与核}

\begin{example}
设 $V$ 是 $\mathbb{R}^n$ 的一个真子空间.判断下列结论是否正确:
\begin{enumerate}
    \item 存在线性变换 $\sigma:\mathbb{R}^n\to\mathbb{R}^n$,使
    \[
    \operatorname{Im}\sigma=V;
    \]
    \item 存在线性变换 $\tau:\mathbb{R}^n\to\mathbb{R}^n$,使
    \[
    \ker\tau=V.
    \]
\end{enumerate}
\end{example}
\exampleblank

\begin{solution}
两者都正确.取 $V$ 的一个补空间 $W$,使
\[
\mathbb{R}^n=V\oplus W.
\]
定义投影
\[
\sigma(v+w)=v,
\]
则 $\operatorname{Im}\sigma=V$.

再取任意与 $W$ 同维的子空间 $U$,以及同构 $T:W\to U$,定义
\[
\tau(v+w)=T(w).
\]
于是
\[
\ker\tau=V.
\]
\end{solution}


\begin{example}
设 $V$ 是复数域上的 $n$ 维线性空间,$\sigma$ 是 $V$ 的线性变换,
$W_1,W_2$ 是任意两个子空间.问
\[
\sigma(W_1\cap W_2)=\sigma(W_1)\cap\sigma(W_2)
\]
是否恒成立?
\end{example}
\exampleblank

\begin{solution}
总有
\[
\sigma(W_1\cap W_2)
\subseteq
\sigma(W_1)\cap\sigma(W_2),
\]
但等号不一定成立.

例如在 $\mathbb{C}^2$ 上取
\[
\sigma(x,y)=(x,0),
\quad
W_1=\Span\{(1,0)\},
\quad
W_2=\Span\{(1,1)\}.
\]
则
\[
W_1\cap W_2=\{0\},
\]
而
\[
\sigma(W_1)=\sigma(W_2)=\Span\{(1,0)\}.
\]
故左端为 $\{0\}$,右端非零.
\end{solution}


\begin{example}
设 $A,B:V\to V$ 为线性变换,$A$ 可逆,$B$ 幂零,且
\[
AB=BA.
\]
证明
\[
\ker(A-B)=\{0\}.
\]
\end{example}
\exampleblank

\begin{solution}
若
\[
(A-B)x=0,
\]
则 $Ax=Bx$.由 $AB=BA$ 归纳得到
\[
A^k x=B^k x.
\]
取 $k$ 使 $B^k=0$,则
\[
A^k x=0.
\]
因 $A^k$ 可逆,故 $x=0$.
\end{solution}


\begin{example}[\markstar]
设 $\sigma$ 是线性空间 $V$ 的线性变换,且
\[
\sigma^2=\sigma.
\]
令
\[
V_1=\operatorname{Im}\sigma,
\qquad
V_2=\ker\sigma.
\]
证明
\[
V=V_1\oplus V_2,
\]
且 $\sigma$ 在 $V_1$ 上为恒等变换.
\end{example}
\exampleblank

\begin{solution}
对任意 $x\in V$,
\[
x=\sigma x+(x-\sigma x).
\]
其中 $\sigma x\in V_1$,并且
\[
\sigma(x-\sigma x)=\sigma x-\sigma^2x=0,
\]
故 $x-\sigma x\in V_2$,从而 $V=V_1+V_2$.

若 $y\in V_1\cap V_2$,则 $y=\sigma x$ 且 $\sigma y=0$,于是
\[
y=\sigma x=\sigma^2x=\sigma y=0.
\]
因此
\[
V=V_1\oplus V_2.
\]
若 $y\in V_1$,写 $y=\sigma x$,则
\[
\sigma y=\sigma^2x=\sigma x=y.
\]
\end{solution}


\begin{example}[\markstar]
设
\[
V=V_1\oplus V_2.
\]
对分解
\[
\alpha=\alpha_1+\alpha_2
\qquad
(\alpha_i\in V_i)
\]
定义
\[
f_1(\alpha)=\alpha_1,\qquad
f_2(\alpha)=\alpha_2.
\]
证明:
\begin{enumerate}
    \item $f_1,f_2$ 是线性变换;
    \item $f_1^2=f_1,\ f_2^2=f_2$;
    \item $f_1f_2=f_2f_1=0$,且 $f_1+f_2=I_V$.
\end{enumerate}
\end{example}
\exampleblank

\begin{solution}
直和分解的唯一性保证 $f_1,f_2$ 定义良好.
对
\[
\alpha=\alpha_1+\alpha_2,\qquad
\beta=\beta_1+\beta_2
\]
直接计算即可得到线性性.又
\[
f_1(\alpha_1)=\alpha_1,\qquad
f_1(\alpha_2)=0,
\]
以及对 $f_2$ 的对应关系,因此
\[
f_i^2=f_i,\qquad
f_1f_2=f_2f_1=0.
\]
最后
\[
(f_1+f_2)(\alpha)=\alpha_1+\alpha_2=\alpha.
\]
\end{solution}


\begin{example}
设 $V_1,V_2$ 是 $n$ 维线性空间 $V$ 的两个子空间,且
\[
\dim V_1+\dim V_2=n.
\]
证明存在 $V$ 的线性变换 $\varphi$ 使
\[
\operatorname{Im}\varphi=V_1,
\qquad
\ker\varphi=V_2.
\]
\end{example}
\exampleblank

\begin{solution}
原题条件本身不足.由秩--零化度定理,若存在这样的 $\varphi$,则
\[
n
=
\dim\operatorname{Im}\varphi+\dim\ker\varphi
=
\dim V_1+\dim V_2,
\]
但还必须有
\[
V_1\cap V_2=\{0\}
\]
才能按投影方式构造.

例如在 $\mathbb{R}^2$ 中取
\[
V_1=V_2=\Span\{(1,0)\},
\]
则维数之和为 $2$,但不存在像空间和核空间都等于同一非零直线的线性变换:
若 $v\in\operatorname{Im}\varphi=\ker\varphi$,则该限制会破坏所需秩结构.

正确命题是:若
\[
V=V_1\oplus V_2,
\]
则取沿 $V_2$ 到 $V_1$ 的投影即可.
\end{solution}


\begin{example}[\markstar]
设 $\sigma$ 是数域 $F$ 上的 $n$ 维线性空间 $V$ 的线性变换,
且
\[
V=V_1\oplus V_2.
\]
证明 $\sigma$ 可逆的充分必要条件是
\[
V=\sigma(V_1)\oplus\sigma(V_2).
\]
\end{example}
\exampleblank

\begin{solution}
若 $\sigma$ 可逆,则
\[
V=\sigma(V)=\sigma(V_1+V_2)=\sigma(V_1)+\sigma(V_2).
\]
若
\[
y\in\sigma(V_1)\cap\sigma(V_2),
\]
有 $y=\sigma(v_1)=\sigma(v_2)$,单射性给出 $v_1=v_2\in V_1\cap V_2$,
故 $y=0$,于是为直和.

反之,若
\[
V=\sigma(V_1)\oplus\sigma(V_2),
\]
则 $\sigma(V)=V$,故 $\sigma$ 满射.有限维情况下满射等价于可逆.
\end{solution}


\begin{proposition}[\marktwostar]
设 $\varphi$ 是 $n$ 维线性空间 $V$ 的线性变换,则下列命题等价:
\begin{enumerate}
    \item $V=\ker\varphi\oplus\operatorname{Im}\varphi$;
    \item $\ker\varphi\cap\operatorname{Im}\varphi=\{0\}$;
    \item $\ker\varphi=\ker\varphi^2$;
    \item $\operatorname{Im}\varphi=\operatorname{Im}\varphi^2$.
\end{enumerate}
\end{proposition}


\begin{example}
设 $\sigma$ 是 $n$ 维线性空间 $V$ 上的线性变换,证明:
\begin{enumerate}
    \item
    \[
    \ker\sigma=\ker\sigma^2
    \iff
    \ker\sigma\cap\operatorname{Im}\sigma=\{0\};
    \]
    \item
    \[
    \operatorname{Im}\sigma=\operatorname{Im}\sigma^2
    \iff
    \ker\sigma+\operatorname{Im}\sigma=V.
    \]
\end{enumerate}
\end{example}
\exampleblank

\begin{solution}
第一问:若 $x=\sigma y\in\ker\sigma\cap\operatorname{Im}\sigma$,则
\[
\sigma^2y=0.
\]
由 $\ker\sigma=\ker\sigma^2$ 得 $\sigma y=0$,即 $x=0$.反向若
$\sigma^2y=0$,则 $\sigma y$ 属于交空间,故 $\sigma y=0$.

第二问使用秩--零化度:
\[
\dim(\ker\sigma+\operatorname{Im}\sigma)
=
\dim\ker\sigma+\dim\operatorname{Im}\sigma
-\dim(\ker\sigma\cap\operatorname{Im}\sigma).
\]
结合第一问以及
\[
\operatorname{Im}\sigma^2\subseteq\operatorname{Im}\sigma
\]
即可得到等价性.
\end{solution}


\begin{example}
设 $\varphi,\psi$ 是线性空间 $V$ 上的幂等变换:
\[
\varphi^2=\varphi,\qquad
\psi^2=\psi.
\]
证明:
\begin{enumerate}
    \item
    \[
    \operatorname{Im}\varphi=\operatorname{Im}\psi
    \iff
    \varphi\psi=\psi,\quad\psi\varphi=\varphi;
    \]
    \item
    \[
    \ker\varphi=\ker\psi
    \iff
    \varphi\psi=\varphi,\quad\psi\varphi=\psi.
    \]
\end{enumerate}
\end{example}
\exampleblank

\begin{solution}
幂等变换在自身像空间上恒等.
若像空间相同,则对任意 $x$,
\[
\psi x\in\operatorname{Im}\psi=\operatorname{Im}\varphi,
\]
故 $\varphi\psi x=\psi x$;另一式同理.反向由
\[
\varphi\psi=\psi
\]
得 $\operatorname{Im}\psi\subseteq\operatorname{Im}\varphi$,
另一式给出反包含.

核空间的证明类似:由
\[
\varphi\psi=\varphi
\]
可得 $\ker\psi\subseteq\ker\varphi$,再由另一式得到反包含.
\end{solution}


\begin{proposition}[\markstar]
设 $V,U,W$ 为有限维线性空间,
\[
\varphi:V\to U,\qquad
\psi:W\to U
\]
为线性映射,则
\[
\exists\,\sigma:V\to W,\quad\varphi=\psi\sigma
\iff
\operatorname{Im}\varphi\subseteq\operatorname{Im}\psi.
\]
\end{proposition}


\begin{example}
设 $f,g$ 是 $n$ 维线性空间 $V$ 的线性变换,且
\[
\ker f\subseteq\ker g.
\]
证明:
\begin{enumerate}
    \item 存在 $V$ 的线性变换 $h$,使 $g=hf$;
    \item 若 $\ker g=\ker f$,则可取 $h$ 为可逆线性变换.
\end{enumerate}
\end{example}
\exampleblank

\begin{solution}
先在 $\operatorname{Im}f$ 上定义
\[
h_0(fx)=gx.
\]
若 $fx=fy$,则 $x-y\in\ker f\subseteq\ker g$,所以 $gx=gy$,
故定义良好.把 $h_0$ 从 $\operatorname{Im}f$ 线性延拓到整个 $V$,
即得到 $g=hf$.

若两核相等,则
\[
\dim\operatorname{Im}f=\dim\operatorname{Im}g,
\]
而 $h_0:\operatorname{Im}f\to\operatorname{Im}g$ 为同构.
分别取补空间并在补空间间任取同构,即可把 $h_0$ 延拓成 $V$ 上的可逆变换.
\end{solution}


\begin{example}
设 $S,T$ 是线性空间 $V$ 的子空间,$f$ 是 $V$ 的线性变换.
记
\[
S^0=\{\ell\in V^*:\ell(S)=0\},
\]
且转置映射
\[
f^*:V^*\to V^*,\qquad
g\mapsto g\circ f.
\]
证明:
\begin{enumerate}
    \item
    \[
    (S\cap T)^0=S^0+T^0;
    \]
    \item
    \[
    \operatorname{Im}f^*=(\ker f)^0.
    \]
\end{enumerate}
\end{example}
\exampleblank

\begin{solution}
第一问一侧包含显然.反向取
\[
\ell\in(S\cap T)^0.
\]
在商空间或分别选取 $S\cap T$ 的补基后,可把 $\ell$ 分解为
\[
\ell=\ell_S+\ell_T,
\qquad
\ell_S\in S^0,\quad \ell_T\in T^0.
\]

第二问若 $g=f^*h=h\circ f$,则 $g$ 在 $\ker f$ 上为零,故
\[
\operatorname{Im}f^*\subseteq(\ker f)^0.
\]
反向若 $g\in(\ker f)^0$,则可在 $\operatorname{Im}f$ 上定义
\[
h_0(fx)=g(x).
\]
这一定义良好,并可延拓成 $V$ 上线性函数 $h$,于是
\[
g=h\circ f=f^*h.
\]
\end{solution}


\begin{theorem}[\marktwostar]
设
\[
f(x)=g(x)h(x),
\qquad
g(x),h(x)\in P[x],
\]
且
\[
(g,h)=1.
\]
若 $\sigma$ 是有限维线性空间 $V$ 上的线性变换,则
\[
\ker f(\sigma)
=
\ker g(\sigma)\oplus\ker h(\sigma).
\]
\end{theorem}

\begin{proof}
由 Bézout 恒等式,存在 $u,v\in P[x]$ 使
\[
u(x)g(x)+v(x)h(x)=1.
\]
对 $\sigma$ 代入:
\[
u(\sigma)g(\sigma)+v(\sigma)h(\sigma)=I.
\]
若 $x\in\ker f(\sigma)$,则
\[
x=u(\sigma)g(\sigma)x+v(\sigma)h(\sigma)x.
\]
第一项属于 $\ker h(\sigma)$,第二项属于 $\ker g(\sigma)$,故核空间为二者之和.
若 $x$ 同时属于两个核,则由 Bézout 式直接得 $x=0$,因此为直和.
\end{proof}


\begin{example}
设 $\sigma$ 是实数域上的 $n$ 维线性空间 $V$ 的线性变换,证明:
\begin{enumerate}
    \item
    \[
    \sigma^3=I
    \iff
    r(I-\sigma)+r(I+\sigma+\sigma^2)=n;
    \]
    \item
    \[
    \sigma^2=I
    \iff
    r(I-\sigma)+r(I+\sigma)=n.
    \]
\end{enumerate}
\end{example}
\exampleblank

\begin{solution}
利用
\[
x^3-1=(x-1)(x^2+x+1),
\qquad
(x-1,x^2+x+1)=1
\]
以及上一条核分解定理.第一式中
\[
\ker(I-\sigma)\oplus\ker(I+\sigma+\sigma^2)=V
\]
恰好等价于秩之和为 $n$,并等价于
\[
(\sigma-I)(\sigma^2+\sigma+I)=\sigma^3-I=0.
\]
第二问同理由
\[
x^2-1=(x-1)(x+1)
\]
得到.
\end{solution}


\begin{example}
设 $A$ 为 $n$ 阶方阵,$B$ 为 $n$ 阶可逆矩阵,且
\[
r(I-AB)+r(I+BA)=n.
\]
证明
\[
r(A)=n.
\]
\end{example}
\exampleblank

\begin{solution}
因为 $B$ 可逆,
\[
BA=B(AB)B^{-1},
\]
故 $AB$ 与 $BA$ 相似,于是
\[
r(I+BA)=r(I+AB).
\]
题设变为
\[
r(I-AB)+r(I+AB)=n.
\]
由上一题可得
\[
(AB)^2=I,
\]
故 $AB$ 可逆.于是 $A$ 也可逆,从而
\[
\boxed{r(A)=n}.
\]
\end{solution}


\begin{example}
设 $A,B$ 是有限维线性空间 $V$ 的线性变换.
\begin{enumerate}
    \item 证明
    \[
    \dim\ker(AB)
    \leq
    \dim\ker A+\dim\ker B;
    \]
    \item 判断
    \[
    \dim\ker(AB)=\dim\ker(BA)
    \]
    是否恒成立.
\end{enumerate}
\end{example}
\exampleblank

\begin{solution}
限制映射
\[
B:\ker(AB)\to\ker A.
\]
其核正是
\[
\ker B.
\]
故秩--零化度给出
\[
\dim\ker(AB)
=
\dim\ker B
+
\dim B(\ker(AB))
\leq
\dim\ker B+\dim\ker A.
\]

第二问不成立.例如
\[
A=
\begin{pmatrix}1&0\\0&0\end{pmatrix},
\qquad
B=
\begin{pmatrix}0&1\\0&0\end{pmatrix}.
\]
则
\[
AB=B,\qquad BA=0,
\]
故
\[
\dim\ker(AB)=1,\qquad
\dim\ker(BA)=2.
\]
\end{solution}


\begin{example}
设 $\varphi_1,\ldots,\varphi_m$ 是 $n$ 维线性空间 $V$ 的线性变换,
满足
\[
\varphi_i^2=\varphi_i,
\qquad
\varphi_i\varphi_j=0\quad(i\neq j).
\]
证明
\[
V
=
\operatorname{Im}\varphi_1\oplus\cdots\oplus\operatorname{Im}\varphi_m
\oplus
\bigcap\limits_{i=1}^{m}\ker\varphi_i.
\]
\end{example}
\exampleblank

\begin{solution}
令
\[
P=\varphi_1+\cdots+\varphi_m.
\]
由两两正交性,
\[
P^2=P.
\]
因此
\[
V=\operatorname{Im}P\oplus\ker P.
\]
又
\[
\operatorname{Im}P
=
\operatorname{Im}\varphi_1\oplus\cdots\oplus\operatorname{Im}\varphi_m,
\]
而
\[
\ker P=\bigcap\limits_{i=1}^{m}\ker\varphi_i.
\]
即得结论.
\end{solution}


\begin{theorem}[\markstar]
设 $A$ 是 $n$ 阶幂等矩阵,即 $A^2=A$.则:
\begin{enumerate}
    \item
    \[
    r(A)+r(A-I)=n;
    \]
    \item $A$ 的特征值只能为 $0$ 或 $1$;
    \item
    \[
    r(A)=\operatorname{tr}A.
    \]
\end{enumerate}
\end{theorem}


\begin{proposition}[\markstar]
设 $A_1,\ldots,A_s$ 为 $n$ 阶矩阵,且
\[
A_1+\cdots+A_s=I_n.
\]
则下列命题等价:
\begin{enumerate}
    \item $A_i$ 全部幂等;
    \item
    \[
    r(A_1)+\cdots+r(A_s)=n;
    \]
    \item
    \[
    A_iA_j=0\qquad(i\neq j).
    \]
\end{enumerate}
\end{proposition}


\begin{example}
设
\[
V=V_1\oplus V_2\oplus\cdots\oplus V_s.
\]
证明存在 $V$ 上的线性变换 $\sigma_1,\ldots,\sigma_s$,满足
\[
\sigma_i^2=\sigma_i,\qquad
\sigma_i\sigma_j=0\ (i\neq j),\qquad
\sigma_1+\cdots+\sigma_s=I,
\]
且
\[
\sigma_i(V)=V_i.
\]
\end{example}
\exampleblank

\begin{solution}
对唯一分解
\[
x=x_1+\cdots+x_s,\qquad x_i\in V_i,
\]
定义
\[
\sigma_i(x)=x_i.
\]
由直和分解的唯一性,各 $\sigma_i$ 线性且满足全部所需关系.
\end{solution}


\begin{example}
设 $A,A_1,A_2$ 为 $n$ 阶方阵,满足
\[
A^2=A,\qquad
A=A_1+A_2,\qquad
r(A)=r(A_1)+r(A_2).
\]
证明
\[
A_i^2=A_i\quad(i=1,2),
\qquad
A_1A_2=A_2A_1=0.
\]
\end{example}
\exampleblank

\begin{solution}
由
\[
\operatorname{Im}A
\subseteq
\operatorname{Im}A_1+\operatorname{Im}A_2
\]
及维数条件可知
\[
\operatorname{Im}A
=
\operatorname{Im}A_1\oplus\operatorname{Im}A_2.
\]
幂等矩阵 $A$ 在 $\operatorname{Im}A$ 上为恒等.若 $y=A_1x$,则
\[
y=Ay=A_1y+A_2y.
\]
按上述直和分解唯一性,
\[
A_1y=y,\qquad A_2y=0.
\]
故
\[
A_1^2=A_1,\qquad A_2A_1=0.
\]
交换 $A_1,A_2$ 即得另外两式.
\end{solution}


\begin{example}
设 $A$ 为幂等矩阵,且
\[
A=A_1+\cdots+A_s,
\qquad
r(A)=r(A_1)+\cdots+r(A_s).
\]
证明所有 $A_i$ 都相似于对角矩阵,且
\[
\operatorname{tr}A_i=r(A_i).
\]
\end{example}
\exampleblank

\begin{solution}
与上一题相同,由秩相加得到
\[
\operatorname{Im}A
=
\operatorname{Im}A_1\oplus\cdots\oplus\operatorname{Im}A_s.
\]
利用 $A$ 在 $\operatorname{Im}A$ 上恒等,对每个分量比较可得
\[
A_i^2=A_i,\qquad A_iA_j=0\quad(i\neq j).
\]
因此每个 $A_i$ 都是幂等矩阵,其极小多项式整除
\[
x(x-1),
\]
无重根,所以可对角化,特征值只有 $0,1$.于是
\[
\boxed{\operatorname{tr}A_i=r(A_i)}.
\]
\end{solution}


\newpage

\section{特征值与特征向量}

\subsection{特征值与特征向量的求解}

\begin{theorem}[\markstar]
设 $S_i$ 是属于互异特征值 $\lambda_i$ 的线性无关特征向量集
$(i=1,\ldots,r)$,则
\[
S_1\cup S_2\cup\cdots\cup S_r
\]
线性无关.特别地,属于不同特征值的非零特征向量线性无关.
\end{theorem}


\begin{proposition}[\markstar]
设
\[
A\in M_{m,n}(K),\qquad
B\in M_{n,m}(K),\qquad m\geq n.
\]
则
\[
|\lambda I_m-AB|
=
\lambda^{m-n}|\lambda I_n-BA|.
\]
因此 $AB$ 与 $BA$ 具有相同的非零特征值,并且代数重数相同.
\end{proposition}


\begin{theorem}[\markstar]
设 $A=(a_{ij})\in M_n(\mathbb{C})$.定义第 $i$ 个 Gershgorin 圆盘
\[
G_i=
\left\{
z\in\mathbb{C}:
|z-a_{ii}|
\leq
\sum\limits_{j\neq i}|a_{ij}|
\right\}.
\]
则 $A$ 的全部特征值都位于
\[
G_1\cup\cdots\cup G_n
\]
中.
\end{theorem}


\begin{theorem}[\markstar]
若 $n$ 个 Gershgorin 圆盘分成若干互不相交的连通区域,其中某一连通区域
恰由 $k$ 个圆盘组成,则该区域内恰有 $k$ 个特征值,重数按代数重数计算.
\end{theorem}


\begin{example}
估计下列矩阵特征值的取值范围,并写出对应 Gershgorin 圆盘:
\[
A_1=
\begin{pmatrix}
1&-0.1&0.2\\
-1.1&-1&0.4\\
-0.3&0.1&0
\end{pmatrix},
\qquad
A_2=
\begin{pmatrix}
-3&1&0.1&-0.5\\
0.1&2&0.4&-0.1\\
1&0.3&0&0\\
0&-2&1&1
\end{pmatrix}.
\]
\end{example}
\exampleblank

\begin{solution}
对 $A_1$,三个圆盘为
\[
|z-1|\leq0.3,\qquad
|z+1|\leq1.5,\qquad
|z|\leq0.4.
\]
故全部特征值都在这三个圆盘之并内.

对 $A_2$,四个圆盘为
\[
|z+3|\leq1.6,\qquad
|z-2|\leq0.6,\qquad
|z|\leq1.3,\qquad
|z-1|\leq3.
\]
由各圆盘的交叠关系还可进一步使用 Gershgorin 第二定理判断各连通分量内
特征值的个数.
\end{solution}


\begin{example}
设 $A$ 为 $n$ 阶方阵.
\begin{enumerate}
    \item 若
    \[
    |a_{ii}|>\sum\limits_{j\neq i}|a_{ij}|
    \qquad(i=1,\ldots,n),
    \]
    证明 $|A|\neq0$;
    \item 若
    \[
    a_{ii}>\sum\limits_{j\neq i}|a_{ij}|
    \qquad(i=1,\ldots,n),
    \]
    证明 $|A|>0$.
\end{enumerate}
\end{example}
\exampleblank

\begin{solution}
第一问中所有 Gershgorin 圆盘都不包含原点,所以 $0$ 不是特征值,
从而 $A$ 可逆.

第二问各圆盘都完全位于右半平面.实矩阵的非实特征值成共轭对出现,
其乘积为正;实特征值全部为正.因此所有特征值之积满足
\[
\boxed{\det A>0}.
\]
原讲义若仅写``严格对角占优''而不要求对角元为正,则第二个结论并不成立;
这里按其前文更精确的条件整理.
\end{solution}


\begin{example}
设 $A=(a_{ij})$ 满足
\[
0\leq a_{ij}\leq1,
\qquad
\sum\limits_{j=1}^{n}a_{ij}=1.
\]
证明:
\begin{enumerate}
    \item $1$ 是 $A$ 的特征值;
    \item 对任意特征值 $\lambda$,都有 $|\lambda|\leq1$.
\end{enumerate}
\end{example}
\exampleblank

\begin{solution}
令
\[
\boldsymbol{1}=(1,\ldots,1)^{\mathsf T}.
\]
由行和为 $1$,
\[
A\boldsymbol{1}=\boldsymbol{1},
\]
故 $1$ 是特征值.

又第 $i$ 个 Gershgorin 圆盘满足
\[
|z-a_{ii}|
\leq
1-a_{ii}.
\]
该圆盘完全包含在单位闭圆盘 $|z|\leq1$ 中,因此所有特征值均满足
\[
|\lambda|\leq1.
\]
\end{solution}


\begin{example}[\markstar]
求
\[
M=
\begin{pmatrix}
0&J_n\\
J_n&0
\end{pmatrix}
\]
的全部特征值以及 $2n$ 个线性无关的特征向量,其中 $J_n$ 为全 $1$ 矩阵.
\end{example}
\exampleblank

\begin{solution}
令
\[
e=(1,\ldots,1)^{\mathsf T}.
\]
则
\[
J_ne=ne,
\]
而对任意 $u\perp e$ 有 $J_nu=0$.

因此
\[
\begin{pmatrix}e\\e\end{pmatrix}
\]
是特征值 $n$ 的特征向量,
\[
\begin{pmatrix}e\\-e\end{pmatrix}
\]
是特征值 $-n$ 的特征向量.

取 $e^\perp$ 的一组基 $u_1,\ldots,u_{n-1}$,则
\[
\begin{pmatrix}u_i\\0\end{pmatrix},
\qquad
\begin{pmatrix}0\\u_i\end{pmatrix}
\]
均属于特征值 $0$.故特征值为
\[
\boxed{n,\ -n,\ 0\ (2n-2\text{ 重})}.
\]
\end{solution}


\begin{example}
设
\[
A=
\begin{pmatrix}
a&-1&c\\
5&b&3\\
1-c&0&-a
\end{pmatrix},
\qquad
|A|=-1.
\]
又 $A^*$ 有特征值 $\lambda_0$,相应特征向量为
\[
\alpha=(-1,-1,1)^{\mathsf T}.
\]
求 $a,b,c,\lambda_0$.
\end{example}
\exampleblank

\begin{solution}
由 $|A|=-1$,矩阵 $A$ 可逆,且
\[
A^*=|A|A^{-1}=-A^{-1}.
\]
条件
\[
A^*\alpha=\lambda_0\alpha
\]
等价于
\[
A\alpha=-\dfrac{1}{\lambda_0}\alpha.
\]
将 $\alpha=(-1,-1,1)^{\mathsf T}$ 代入并结合 $\det A=-1$ 解方程组,
得到
\[
\boxed{
a=2,\qquad b=-3,\qquad c=2,\qquad \lambda_0=1.
}
\]
\end{solution}


\begin{example}
已知三阶方阵 $A$ 满足
\[
|A-I|=|A-2I|=|A+I|=\lambda.
\]
\begin{enumerate}
    \item 当 $\lambda=0$ 时,求 $|A+3I|$;
    \item 当 $\lambda=2$ 时,求 $|A+3I|$.
\end{enumerate}
\end{example}
\exampleblank

\begin{solution}
令
\[
p(t)=\det(A-tI).
\]
这是首项系数为 $-1$ 的三次多项式,并且
\[
p(1)=p(2)=p(-1)=\lambda.
\]
故
\[
p(t)-\lambda=-(t-1)(t-2)(t+1).
\]
于是
\[
|A+3I|=p(-3)=\lambda+40.
\]
因此
\[
\boxed{
\lambda=0\Rightarrow |A+3I|=40,
\qquad
\lambda=2\Rightarrow |A+3I|=42.
}
\]
\end{solution}


\begin{example}
已知 $n$ 阶方阵
\[
A=(a_ia_j)_{n\times n}+J_n-I_n,
\]
即
\[
a_{ii}=a_i^2,\qquad
a_{ij}=a_ia_j+1\quad(i\neq j),
\]
并且
\[
\sum\limits_{i=1}^{n}a_i=1,
\qquad
\sum\limits_{i=1}^{n}a_i^2=n.
\]
求:
\begin{enumerate}
    \item $A$ 的全部特征值;
    \item $\det A$ 与 $\operatorname{tr}A$.
\end{enumerate}
\end{example}
\exampleblank

\begin{solution}
令
\[
a=(a_1,\ldots,a_n)^{\mathsf T},
\qquad
e=(1,\ldots,1)^{\mathsf T}.
\]
则
\[
A=aa^{\mathsf T}+ee^{\mathsf T}-I.
\]
且
\[
a^{\mathsf T}a=n,\qquad
e^{\mathsf T}e=n,\qquad
a^{\mathsf T}e=1.
\]
在
\[
\Span\{a,e\}
\]
上,相对于基 $(a,e)$,$A$ 的矩阵为
\[
\begin{pmatrix}
n-1&1\\
1&n-1
\end{pmatrix},
\]
故得到特征值
\[
n,\qquad n-2.
\]
在同时正交于 $a,e$ 的 $n-2$ 维子空间上,
\[
A=-I.
\]
因此
\[
\boxed{
\lambda(A):
n,\ n-2,\ \underbrace{-1,\ldots,-1}_{n-2}
}.
\]
从而
\[
\boxed{
\det A=(-1)^{n-2}n(n-2),
\qquad
\operatorname{tr}A=n.
}
\]
原 PDF 中第一条件印成了 $\sum a_1=1$,按上下文应为 $\sum a_i=1$.
\end{solution}


\begin{example}[\markstar]
设 $n$ 阶循环矩阵
\[
C=
\begin{pmatrix}
c_0&c_1&\cdots&c_{n-1}\\
c_{n-1}&c_0&\cdots&c_{n-2}\\
\vdots&\vdots&&\vdots\\
c_1&c_2&\cdots&c_0
\end{pmatrix}.
\]
\begin{enumerate}
    \item 求全部特征值及相应特征向量;
    \item 求 $\det C$.
\end{enumerate}
\end{example}
\exampleblank

\begin{solution}
令
\[
f(x)=c_0+c_1x+\cdots+c_{n-1}x^{n-1},
\]
并令 $\omega_k$ 遍历全部 $n$ 次单位根.则
\[
v_k=
(1,\omega_k,\omega_k^2,\ldots,\omega_k^{n-1})^{\mathsf T}
\]
是特征向量,对应特征值
\[
\boxed{\lambda_k=f(\omega_k)}.
\]
因此
\[
\boxed{
\det C
=
\prod\limits_{k=1}^{n}f(\omega_k).
}
\]
\end{solution}


\begin{example}
设 $A$ 是数域 $K$ 上的 $n$ 阶方阵,$\lambda_0\in K$ 是特征多项式的单根,
$\alpha$ 是属于 $\lambda_0$ 的特征向量.证明
\[
(\lambda_0I-A)x=\alpha
\]
无解.
\end{example}
\exampleblank

\begin{solution}
若存在解 $x$,则
\[
(A-\lambda_0I)x=-\alpha,
\qquad
(A-\lambda_0I)\alpha=0.
\]
于是
\[
(A-\lambda_0I)^2x=0,
\qquad
(A-\lambda_0I)x\neq0.
\]
这说明 $\lambda_0$ 至少存在长度为 $2$ 的广义特征向量链,因此其代数重数至少为 $2$,
与 $\lambda_0$ 为单根矛盾.
\end{solution}


\begin{example}
设 $A$ 是 $n$ 阶复方阵,且存在列向量 $\alpha$,使
\[
\alpha,A\alpha,\ldots,A^{n-1}\alpha
\]
线性无关.证明对任意特征值 $\lambda_0$,
其特征子空间维数为
\[
\dim V_{\lambda_0}=1.
\]
\end{example}
\exampleblank

\begin{solution}
给定的 $n$ 个向量构成一组基,因此 $A$ 在这组循环基下的矩阵是一个友矩阵.
友矩阵的每个特征值对应的特征子空间维数均为 $1$:
解
\[
(F-\lambda I)x=0
\]
时前 $n-1$ 个分量都由一个自由参数唯一确定.
相似变换保持特征子空间维数,故
\[
\boxed{\dim V_{\lambda_0}=1}.
\]
\end{solution}


\begin{example}
设 $A$ 为 $n$ 阶不可逆矩阵.证明 $A^*$ 的 $n$ 个特征值中至少有
$n-1$ 个为 $0$,且若存在非零特征值,则它等于
\[
\operatorname{tr}A^*.
\]
\end{example}
\exampleblank

\begin{solution}
若
\[
r(A)\leq n-2,
\]
则所有 $n-1$ 阶子式均为零,所以
\[
A^*=0.
\]
若
\[
r(A)=n-1,
\]
则
\[
r(A^*)=1.
\]
秩一矩阵的特征值为
\[
0\quad(n-1\text{ 重})
\]
以及
\[
\operatorname{tr}A^*
\]
一次.因此结论成立.
\end{solution}


\begin{example}
设 $A$ 是数域 $K$ 上的 $n$ 阶可逆矩阵,并且
\[
A\sim A^k
\qquad
(\forall k\in\mathbb{Z}_+).
\]
证明 $A$ 的全部特征值都是 $1$.
\end{example}
\exampleblank

\begin{solution}
设谱的有限集合为
\[
S=\{\lambda_1,\ldots,\lambda_s\}\subset\overline K.
\]
由
\[
A\sim A^k
\]
可知对每个正整数 $k$,
\[
S=S^k=\{\lambda_1^k,\ldots,\lambda_s^k\}.
\]
因此每个 $\lambda_i$ 的所有正整数幂仍落在有限集合 $S$ 中,所以 $\lambda_i$ 必为单位根.
取 $N$ 为这些单位根阶数的公倍数,则
\[
A^N
\]
的全部特征值都为 $1$.又 $A\sim A^N$,故 $A$ 本身的全部特征值也都为 $1$.
\end{solution}


\begin{example}
设 $a_{ij}\in\mathbb{Z}$,证明
\[
D_n=
\det
\begin{pmatrix}
a_{11}-\dfrac13&a_{12}&\cdots&a_{1n}\\
a_{21}&a_{22}-\dfrac13&\cdots&a_{2n}\\
\vdots&\vdots&\ddots&\vdots\\
a_{n1}&a_{n2}&\cdots&a_{nn}-\dfrac13
\end{pmatrix}
\neq0.
\]
\end{example}
\exampleblank

\begin{solution}
若 $D_n=0$,则 $\dfrac13$ 是整数矩阵
\[
A=(a_{ij})
\]
的特征值.于是 $\dfrac13$ 是其首一整系数特征多项式的有理根.
由有理根定理,首一整系数多项式的有理根必为整数,矛盾.
故
\[
\boxed{D_n\neq0}.
\]
\end{solution}


\subsection{对角化}

\begin{theorem}[\markstar]
若 $n$ 阶复矩阵 $A$ 有 $n$ 个互不相同的特征值,则 $A$ 可对角化.
\end{theorem}


\begin{proposition}[\marktwostar]
设 $A$ 为 $n$ 阶复矩阵,则下列命题等价:
\begin{enumerate}
    \item $A$ 可对角化;
    \item $A$ 有 $n$ 个线性无关的特征向量;
    \item $\mathbb{C}^n$ 是各特征子空间的直和;
    \item 每个特征值的代数重数等于几何重数;
    \item 对任意 $A$-不变子空间 $U$,存在 $A$-不变子空间 $W$ 使
    \[
    \mathbb{C}^n=U\oplus W;
    \]
    \item $A$ 的极小多项式无重根;
    \item $A$ 的所有 Jordan 块都是一阶的.
\end{enumerate}
\end{proposition}


\begin{example}
已知
\[
A=
\begin{pmatrix}
1&2&2\\
2&a&2\\
2&2&1
\end{pmatrix},
\qquad
B=
\begin{pmatrix}
-1&0&0\\
0&-1&0\\
0&0&b
\end{pmatrix}.
\]
求 $a,b$ 使 $A\sim B$,并求可逆矩阵 $P$ 使
\[
P^{-1}AP=B.
\]
\end{example}
\exampleblank

\begin{solution}
计算
\[
\chi_A(\lambda)
=
(\lambda+1)
\left(
\lambda^2-(a+3)\lambda+3a-8
\right).
\]
要与
\[
(\lambda+1)^2(\lambda-b)
\]
相同,必须
\[
a=1,\qquad b=5.
\]
当 $a=1$ 时,属于 $-1$ 的两个线性无关特征向量可取
\[
(-1,1,0)^{\mathsf T},
\qquad
(-1,0,1)^{\mathsf T},
\]
属于 $5$ 的特征向量可取
\[
(1,1,1)^{\mathsf T}.
\]
因此
\[
\boxed{
P=
\begin{pmatrix}
-1&-1&1\\
1&0&1\\
0&1&1
\end{pmatrix}
}
\]
满足
\[
P^{-1}AP=\diag(-1,-1,5).
\]
\end{solution}


\begin{example}[\markstar]
设 $A,B$ 为 $n$ 阶实方阵,且存在复可逆矩阵 $Q$ 使
\[
A=Q^{-1}BQ.
\]
证明存在实可逆矩阵 $P$ 使
\[
A=P^{-1}BP.
\]
\end{example}
\exampleblank

\begin{solution}
写
\[
Q=X+iY,
\qquad
X,Y\in M_n(\mathbb{R}).
\]
由
\[
BQ=QA
\]
比较实部,虚部得
\[
BX=XA,\qquad BY=YA.
\]
故对任意实数 $t$,
\[
B(X+tY)=(X+tY)A.
\]
多项式
\[
p(t)=\det(X+tY)
\]
不恒为零,否则在复数 $t=i$ 处也应为零,与 $\det Q\neq0$ 矛盾.
因此可取某个实数 $t_0$ 使
\[
P=X+t_0Y
\]
可逆,于是
\[
BP=PA,
\qquad
A=P^{-1}BP.
\]
\end{solution}


\begin{example}
设三阶方阵 $A$ 的特征值为
\[
1,\quad2,\quad3,
\]
对应特征向量依次为
\[
\xi_1=(1,1,1)^{\mathsf T},\quad
\xi_2=(1,2,4)^{\mathsf T},\quad
\xi_3=(1,3,9)^{\mathsf T}.
\]
又
\[
\beta=(1,1,3)^{\mathsf T}.
\]
\begin{enumerate}
    \item 将 $\beta$ 用 $\xi_1,\xi_2,\xi_3$ 线性表示;
    \item 求 $A^n\beta$.
\end{enumerate}
\end{example}
\exampleblank

\begin{solution}
解
\[
\beta=c_1\xi_1+c_2\xi_2+c_3\xi_3
\]
得
\[
(c_1,c_2,c_3)=(2,-2,1).
\]
因此
\[
\boxed{
\beta=2\xi_1-2\xi_2+\xi_3
}.
\]
利用特征向量关系,
\[
\boxed{
A^n\beta
=
2\xi_1-2\cdot2^n\xi_2+3^n\xi_3.
}
\]
\end{solution}


\begin{example}
用 $J$ 表示全 $1$ 的 $n$ 阶方阵,$n\geq2$,
\[
f(x)=a+bx\in\mathbb{Q}[x],
\qquad
A=f(J)=aI+bJ.
\]
\begin{enumerate}
    \item 求 $J$ 的全部特征值和特征向量;
    \item 求 $A$ 的全部特征子空间;
    \item 判断 $A$ 是否可对角化,并给出一个有理可逆矩阵把它对角化.
\end{enumerate}
\end{example}
\exampleblank

\begin{solution}
令
\[
e=(1,\ldots,1)^{\mathsf T}.
\]
则
\[
Je=ne,
\]
而
\[
e^\perp=\left\{x:\sum\limits_i x_i=0\right\}
\]
是特征值 $0$ 的特征子空间,维数 $n-1$.

故 $A=aI+bJ$ 的特征值为
\[
a+bn
\]
在 $\Span\{e\}$ 上,以及
\[
a
\]
在 $e^\perp$ 上.若 $b\neq0$,这两个特征子空间即全部;若 $b=0$,则
$A=aI$,整个空间都是特征子空间.

取
\[
e,\quad e_1-e_n,\quad e_2-e_n,\ldots,e_{n-1}-e_n
\]
为基即可得到有理可逆对角化矩阵,因此 $A$ 总可对角化.
\end{solution}


\begin{example}
设三阶实对称矩阵 $A$ 的各行元素之和均为 $3$,
\[
\alpha_1=(-1,2,-1)^{\mathsf T},
\qquad
\alpha_2=(0,-1,1)^{\mathsf T}
\]
是 $Ax=0$ 的两个解.
\begin{enumerate}
    \item 求 $A$ 的特征值与特征向量;
    \item 求正交矩阵 $Q$ 和对角矩阵 $D$,使
    \[
    Q^{\mathsf T}AQ=D;
    \]
    \item 求
    \[
    \left|
    \left(\dfrac23B^2\right)^{-1}
    +\dfrac49B^*
    +B
    \right|,
    \]
    其中 $B$ 与 $A-\dfrac32I$ 相似.
\end{enumerate}
\end{example}
\exampleblank

\begin{solution}
由于行和为 $3$,
\[
e=(1,1,1)^{\mathsf T}
\]
是特征值 $3$ 的特征向量.又 $\alpha_1,\alpha_2$ 线性无关并属于核,
故特征值为
\[
\boxed{0,0,3}.
\]
三者互相正交化,单位化即可得到正交矩阵 $Q$,并有
\[
D=\diag(0,0,3).
\]
事实上这也说明
\[
A=J_3.
\]

于是 $B$ 的特征值为
\[
-\dfrac32,\quad-\dfrac32,\quad\dfrac32.
\]
从而
\[
B^2=\dfrac94I,
\qquad
\left(\dfrac23B^2\right)^{-1}=\dfrac23I.
\]
又
\[
\det B=\dfrac{27}{8},
\]
所以 $B^*$ 在对应特征方向上的特征值为
\[
-\dfrac94,\quad-\dfrac94,\quad\dfrac94.
\]
因此题中矩阵的三个特征值为
\[
-\dfrac{11}{6},\quad
-\dfrac{11}{6},\quad
\dfrac{19}{6},
\]
故
\[
\boxed{
\det
=
\dfrac{2299}{216}.
}
\]
\end{solution}


\begin{example}
设 $A\in M_n(K)$,证明:若
\[
r(A)+r(A-I)=n,
\]
则 $A$ 可对角化.
\end{example}
\exampleblank

\begin{solution}
由秩--零化度,
\[
\dim\ker A+\dim\ker(A-I)
=
2n-r(A)-r(A-I)=n.
\]
又
\[
\ker A\cap\ker(A-I)=\{0\}.
\]
故
\[
K^n=\ker A\oplus\ker(A-I).
\]
这两个子空间分别是特征值 $0,1$ 的特征子空间,因此可取由特征向量组成的基,
$A$ 可对角化.
\end{solution}


\begin{example}
设
\[
V=M_n(F),
\qquad
\sigma(A)=A^{\mathsf T}.
\]
\begin{enumerate}
    \item 证明 $\sigma$ 是线性变换;
    \item 求全部特征子空间;
    \item 证明 $\sigma$ 可对角化.
\end{enumerate}
\end{example}
\exampleblank

\begin{solution}
线性性由转置的线性性直接得到.
若底域特征不为 $2$,则
\[
\sigma(A)=A
\]
对应特征值 $1$,特征子空间为全体对称矩阵;
\[
\sigma(A)=-A
\]
对应特征值 $-1$,特征子空间为全体反对称矩阵.
任意矩阵有唯一分解
\[
A=
\dfrac{A+A^{\mathsf T}}{2}
+
\dfrac{A-A^{\mathsf T}}{2},
\]
故
\[
M_n(F)
=
\operatorname{Sym}_n(F)
\oplus
\operatorname{Skew}_n(F),
\]
所以 $\sigma$ 可对角化.
\end{solution}


\begin{example}
设 $A^2=I$.证明对任意正整数 $m,k$,
\[
r((A+I)^m)+r((A-I)^k)=n.
\]
\end{example}
\exampleblank

\begin{solution}
在特征不为 $2$ 的数域上,
\[
x^2-1=(x-1)(x+1)
\]
无重根,所以
\[
V=\ker(A-I)\oplus\ker(A+I).
\]
在 $\ker(A-I)$ 上,
\[
A+I=2I,\qquad A-I=0;
\]
在 $\ker(A+I)$ 上恰好相反.因此
\[
r((A+I)^m)=\dim\ker(A-I),
\]
\[
r((A-I)^k)=\dim\ker(A+I),
\]
二者之和为 $n$.
\end{solution}


\subsection{多个矩阵同时对角化}

\begin{proposition}[\markstar]
设 $A,B\in M_n(K)$ 且
\[
AB=BA.
\]
若 $A,B$ 都可对角化,则它们可同时对角化.
\end{proposition}


\begin{example}[\markstar]
设 $S$ 是任意多个可相似对角化的 $n$ 阶方阵构成的集合,且其中任意两矩阵可交换.
证明 $S$ 中全部矩阵可同时相似对角化.
\end{example}
\exampleblank

\begin{solution}
对维数归纳.若 $S$ 中所有矩阵都是数量矩阵则显然成立.
否则取一个非数量矩阵 $A\in S$.因 $A$ 可对角化,
\[
V=V_{\lambda_1}\oplus\cdots\oplus V_{\lambda_r},
\qquad r\geq2.
\]
任意 $B\in S$ 与 $A$ 可交换,因此保持每个 $V_{\lambda_i}$ 不变.
所有 $B$ 在每个较低维子空间上的限制仍两两可交换且可对角化.
由归纳假设,在每个 $V_{\lambda_i}$ 上存在公共特征基.
合并这些基即得到整个 $S$ 的公共特征基.
\end{solution}


\begin{example}
设 $A$ 有 $n$ 个互不相同的特征值.若 $B$ 满足
\[
AB=BA,
\]
证明 $B$ 可对角化,且 $B$ 可以表示为 $A$ 的多项式.
\end{example}
\exampleblank

\begin{solution}
取 $A$ 的特征基
\[
v_1,\ldots,v_n,
\qquad
Av_i=\lambda_i v_i.
\]
由 $AB=BA$,
\[
A(Bv_i)=B(Av_i)=\lambda_iBv_i,
\]
故 $Bv_i$ 仍属于一维特征子空间 $\Span\{v_i\}$,于是
\[
Bv_i=\mu_i v_i.
\]
因此 $A,B$ 在同一基下均为对角矩阵.

由于 $\lambda_i$ 两两不同,由 Lagrange 插值存在次数小于 $n$ 的多项式 $p$ 满足
\[
p(\lambda_i)=\mu_i.
\]
于是
\[
\boxed{B=p(A)}.
\]
\end{solution}


\begin{example}
多项式
\[
f(x)=x^n+c_{n-1}x^{n-1}+\cdots+c_1x+c_0
\]
的友矩阵为
\[
F=
\begin{pmatrix}
0&0&\cdots&0&-c_0\\
1&0&\cdots&0&-c_1\\
0&1&\cdots&0&-c_2\\
\vdots&\vdots&\ddots&\vdots&\vdots\\
0&0&\cdots&1&-c_{n-1}
\end{pmatrix}.
\]
证明:若方阵 $A$ 相似于某个多项式的友矩阵,则所有与 $A$ 可交换的方阵都是 $A$ 的多项式.
\end{example}
\exampleblank

\begin{solution}
友矩阵 $F$ 是循环矩阵意义下的循环算子:
取第一标准基向量 $e_1$,则
\[
e_1,Fe_1,\ldots,F^{n-1}e_1
\]
构成一组基.

若 $B$ 与 $F$ 可交换,写
\[
Be_1=
c_0e_1+c_1Fe_1+\cdots+c_{n-1}F^{n-1}e_1
=
p(F)e_1.
\]
于是对任意 $k$,
\[
BF^ke_1
=
F^kBe_1
=
p(F)F^ke_1.
\]
因为这些向量构成基,故
\[
B=p(F).
\]
若 $A=P^{-1}FP$,则与 $A$ 可交换的矩阵共轭到与 $F$ 可交换的矩阵,
所以同样是 $A$ 的多项式.
\end{solution}


\newpage

\section{不变子空间}

\begin{example}
已知线性空间 $V$ 的线性变换 $\sigma$ 在基
$\eta_1,\eta_2,\eta_3,\eta_4$ 下的矩阵为
\[
A=
\begin{pmatrix}
1&-1&-1&2\\
0&1&0&0\\
2&3&1&-1\\
1&-2&-2&-1
\end{pmatrix}.
\]
求包含 $\eta_1$ 的最小 $\sigma$-不变子空间.
\end{example}
\exampleblank

\begin{solution}
包含 $\eta_1$ 的最小不变子空间就是循环子空间
\[
W=\Span\{\eta_1,\sigma\eta_1,\sigma^2\eta_1,\ldots\}.
\]
由矩阵第一列,
\[
[\sigma\eta_1]_{\eta}
=
\begin{pmatrix}1\\0\\2\\1\end{pmatrix},
\]
而
\[
[\sigma^2\eta_1]_{\eta}
=
\begin{pmatrix}1\\0\\3\\-4\end{pmatrix}.
\]
这三个向量线性无关,而再作用一次 $\sigma$ 后不再增加秩.因此
\[
\boxed{
W=
\Span\left\{
\eta_1,\,
\eta_1+2\eta_3+\eta_4,\,
\eta_1+3\eta_3-4\eta_4
\right\}.
}
\]
\end{solution}


\begin{example}[\markstar]
设 $\sigma$ 是实数域上的 $n$ 维线性空间 $V$ 的线性变换,
$\alpha\neq0$,令
\[
W=\Span\{\alpha,\sigma\alpha,\sigma^2\alpha,\ldots\}.
\]
\begin{enumerate}
    \item 证明 $W$ 是 $\sigma$-不变子空间;
    \item 若 $\dim W=r$,证明
    \[
    \alpha,\sigma\alpha,\ldots,\sigma^{r-1}\alpha
    \]
    是 $W$ 的一组基,并求 $\sigma|_W$ 在此基下的矩阵.
\end{enumerate}
\end{example}
\exampleblank

\begin{solution}
第一问由
\[
\sigma(\sigma^k\alpha)=\sigma^{k+1}\alpha
\]
立即得到.

设
\[
\alpha,\sigma\alpha,\ldots,\sigma^{r-1}\alpha
\]
在最早位置出现线性相关.若在第 $s<r$ 项之前已经相关,则所有更高次幂都可递推表示,
从而 $\dim W\leq s<r$,矛盾.因此前 $r$ 个向量线性无关并构成基.

存在唯一的 $c_0,\ldots,c_{r-1}$ 使
\[
\sigma^r\alpha
=
c_0\alpha+c_1\sigma\alpha+\cdots+c_{r-1}\sigma^{r-1}\alpha.
\]
故限制变换的矩阵为友矩阵型
\[
\boxed{
\begin{pmatrix}
0&0&\cdots&0&c_0\\
1&0&\cdots&0&c_1\\
0&1&\ddots&\vdots&c_2\\
\vdots&\ddots&\ddots&0&\vdots\\
0&\cdots&0&1&c_{r-1}
\end{pmatrix}.
}
\]
\end{solution}


\begin{example}
设 $\varphi$ 是 $n$ 维线性空间 $V$ 的可逆线性变换,
$W$ 是 $\varphi$-不变子空间.证明 $W$ 也是 $\varphi^{-1}$-不变子空间.
\end{example}
\exampleblank

\begin{solution}
由不变性,
\[
\varphi(W)\subseteq W.
\]
因为 $\varphi$ 可逆,所以 $\varphi|_W$ 是单射,从而
\[
\dim\varphi(W)=\dim W.
\]
有限维下包含且维数相同意味着
\[
\varphi(W)=W.
\]
两边作用 $\varphi^{-1}$ 即得
\[
\varphi^{-1}(W)=W.
\]
\end{solution}


\begin{example}[\markstar]
设 $A,B$ 都是线性空间 $V$ 的线性变换,并且
\[
AB=BA.
\]
证明:
\begin{enumerate}
    \item $A$ 的每个特征子空间都是 $B$ 的不变子空间;
    \item 若 $V$ 是复线性空间,则 $A,B$ 至少有一个公共特征向量.
\end{enumerate}
\end{example}
\exampleblank

\begin{solution}
若
\[
x\in V_\lambda(A),
\]
则
\[
A(Bx)=B(Ax)=B(\lambda x)=\lambda Bx,
\]
故
\[
Bx\in V_\lambda(A).
\]
第一问成立.

在复数域上,$A$ 至少有一个特征值 $\lambda$.
若相应特征子空间 $V_\lambda(A)$ 非零,则由第一问它对 $B$ 不变.
限制变换
\[
B|_{V_\lambda(A)}
\]
仍是复线性变换,因此有特征向量 $v\neq0$.这个 $v$ 同时是 $A$ 与 $B$ 的特征向量.
\end{solution}


\begin{example}[\markstar]
设 $\{\sigma_i:i\in I\}$ 是复数域上的 $n$ 维线性空间 $V$ 上的一族
两两可交换线性变换,指标集 $I$ 可以有限也可以无限.证明所有 $\sigma_i$
至少有一个公共特征向量.
\end{example}
\exampleblank

\begin{solution}
对 $\dim V$ 归纳.若所有 $\sigma_i$ 都是数量变换,则任意非零向量都是公共特征向量.

否则取一个非数量变换 $\sigma_{i_0}$,选择其某个非零真特征子空间
\[
W=V_\lambda(\sigma_{i_0}).
\]
由于全部变换两两可交换,上一题表明 $W$ 对每个 $\sigma_i$ 都不变.
于是所有限制
\[
\sigma_i|_W
\]
仍两两可交换,而 $\dim W<\dim V$.由归纳假设,$W$ 中存在公共特征向量,
它也是原来全部 $\sigma_i$ 的公共特征向量.
\end{solution}


\begin{example}[\markstar]
设 $\sigma$ 是 $n>1$ 维线性空间 $V$ 的线性变换,在某一基下的矩阵为单个
Jordan 块
\[
J_n(\lambda)=
\begin{pmatrix}
\lambda&1&0&\cdots&0\\
0&\lambda&1&\ddots&\vdots\\
\vdots&\ddots&\ddots&\ddots&0\\
0&\cdots&0&\lambda&1\\
0&\cdots&\cdots&0&\lambda
\end{pmatrix}.
\]
\begin{enumerate}
    \item 证明 $V$ 不能分解成两个非平凡 $\sigma$-不变子空间的直和;
    \item 求 $\sigma$ 的全部不变子空间.
\end{enumerate}
\end{example}
\exampleblank

\begin{solution}
令
\[
N=\sigma-\lambda I.
\]
则 $N$ 是单个 $n$ 阶幂零 Jordan 块.不变子空间对 $\sigma$ 不变当且仅当对
$N$ 不变.

设对应基为 $e_1,\ldots,e_n$,满足
\[
Ne_1=0,\qquad
Ne_j=e_{j-1}\quad(j\geq2).
\]
若非零不变子空间 $W$ 中某向量的最高非零坐标出现在第 $k$ 位,
连续作用 $N$ 可推出
\[
e_1,\ldots,e_k\in W.
\]
再由最高位的最大性得到
\[
W=\Span\{e_1,\ldots,e_k\}.
\]
因此全部不变子空间恰为
\[
\boxed{
0,\ 
\Span\{e_1\},\
\Span\{e_1,e_2\},\ldots,
\Span\{e_1,\ldots,e_n\}=V.
}
\]
这些子空间构成一条严格嵌套链,所以两个非平凡不变子空间不可能互补成直和.
\end{solution}


\begin{example}[\markstar]
设 $V$ 是数域 $P$ 上的 $n$ 维线性空间,$\sigma$ 满足
\[
\sigma^{n-1}\neq0,\qquad
\sigma^n=0.
\]
证明 $\sigma$ 恰有 $n+1$ 个不变子空间.
\end{example}
\exampleblank

\begin{solution}
由 $\sigma^{n-1}\neq0$,存在 $v$ 使
\[
\sigma^{n-1}v\neq0.
\]
于是
\[
\sigma^{n-1}v,\sigma^{n-2}v,\ldots,\sigma v,v
\]
构成 $V$ 的一组基,$\sigma$ 在此基下就是单个 $n$ 阶幂零 Jordan 块.
由上一题,全部不变子空间恰好为
\[
0,\ 
\Span\{\sigma^{n-1}v\},\
\Span\{\sigma^{n-1}v,\sigma^{n-2}v\},
\ldots,V,
\]
共
\[
\boxed{n+1}
\]
个.
\end{solution}


\begin{example}
设 $\varphi$ 是复线性空间 $V$ 的线性变换,
$\lambda_1,\ldots,\lambda_k$ 是互不相同的特征值,
$v_i$ 是属于 $\lambda_i$ 的特征向量.
设 $W$ 是 $\varphi$-不变子空间,且
\[
w=c_1v_1+\cdots+c_kv_k\in W,
\qquad
c_i\neq0.
\]
证明
\[
v_i\in W\qquad(i=1,\ldots,k).
\]
\end{example}
\exampleblank

\begin{solution}
对每个 $i$,取 Lagrange 插值多项式
\[
p_i(x)
=
\prod_{\substack{1\leq j\leq k\\j\neq i}}
\dfrac{x-\lambda_j}{\lambda_i-\lambda_j}.
\]
则
\[
p_i(\lambda_j)=\delta_{ij}.
\]
因为 $W$ 对 $\varphi$ 不变,也对任意多项式 $p_i(\varphi)$ 不变,所以
\[
p_i(\varphi)w
=
c_i v_i
\in W.
\]
因 $c_i\neq0$,故
\[
v_i\in W.
\]
\end{solution}


\begin{example}[\markstar]
设 $V$ 是实数域上的 $n$ 维线性空间,$\sigma$ 是 $V$ 上的线性变换.
证明 $\sigma$ 必有一维或二维不变子空间.
\end{example}
\exampleblank

\begin{solution}
$\sigma$ 的实特征多项式次数为 $n$.若它有实根 $\lambda$,则对应实特征向量张成一维不变子空间.

若没有实特征值,则必有一对共轭非实特征值
\[
a\pm bi,\qquad b\neq0.
\]
取复特征向量
\[
z=u+iv.
\]
比较实部,虚部可得
\[
\sigma u=au-bv,\qquad
\sigma v=bu+av.
\]
因此
\[
\Span_{\mathbb{R}}\{u,v\}
\]
是二维实不变子空间.
\end{solution}


\begin{example}
设 $V$ 是复数域上的 $n$ 维线性空间,$\varphi$ 是 $V$ 的线性变换.
证明存在不变子空间链
\[
V_1\subseteq V_2\subseteq\cdots\subseteq V_n=V,
\qquad
\dim V_i=i.
\]
\end{example}
\exampleblank

\begin{solution}
复数域上 $\varphi$ 至少有一个特征向量,令其张成 $V_1$.
在商空间
\[
V/V_1
\]
上诱导线性变换,再取一个一维不变子空间并取其逆像得到 $V_2$.
如此递归即可构造
\[
V_1\subseteq\cdots\subseteq V_n.
\]
等价地,这正是``复矩阵可上三角化''的几何表述.
\end{solution}


\begin{example}[\markstar]
设 $\sigma$ 是数域 $P$ 上 $n$ 维线性空间 $V$ 的线性变换,
且有 $n$ 个互不相同的特征值.证明 $V$ 中恰有
\[
2^n
\]
个 $\sigma$-不变子空间.
\end{example}
\exampleblank

\begin{solution}
取对应特征向量基
\[
v_1,\ldots,v_n.
\]
任取指标子集
\[
S\subseteq\{1,\ldots,n\},
\]
子空间
\[
W_S=\Span\{v_i:i\in S\}
\]
显然 $\sigma$-不变,因此至少有 $2^n$ 个.

反之,若 $W$ 是任意不变子空间,对其中任意向量
\[
w=\sum c_i v_i
\]
应用上一题的 Lagrange 投影可知,只要 $c_i\neq0$,就有 $v_i\in W$.
因此 $W$ 必恰好是某些 $v_i$ 的张成空间.故总数正好为
\[
\boxed{2^n}.
\]
\end{solution}


\begin{example}[\markstar]
设 $V$ 是复数域上的有限维线性空间,$\tau$ 是 $V$ 的线性变换.
若对任意 $\tau$-不变子空间 $U$,都存在 $\tau$-不变子空间 $W$ 使
\[
V=U\oplus W,
\]
则称 $\tau$ 完全可约.证明:$\tau$ 完全可约当且仅当 $V$ 有一组由特征向量组成的基.
\end{example}
\exampleblank

\begin{solution}
若 $V$ 有特征向量基,则 $\tau$ 可对角化.对任意不变子空间 $U$,
可由不同特征值的谱投影把 $U$ 分解为各特征子空间中的部分,再在每个特征子空间内取补空间,
这些补空间的直和给出 $\tau$-不变补空间.

反之,设 $\tau$ 完全可约.复数域上存在特征向量 $v_1$,令
\[
U_1=\Span\{v_1\}.
\]
完全可约性给出不变补空间
\[
V=U_1\oplus W_1.
\]
对 $W_1$ 上的限制重复这一过程.有限维归纳最终得到
\[
V
=
\Span\{v_1\}\oplus\cdots\oplus\Span\{v_n\},
\]
其中每个 $v_i$ 都是 $\tau$ 的特征向量.
\end{solution}
